\section{Checking the labels in real pods}
\label{sec:validation}

The Job labels in \cref{sec:accounting} are derived from the workload's measurement, so we
checked them twice. The first check (rt-k8sval-2) re-ran 176 corpus runs both as Flyte
tasks and as their rendered \code{job.py}, as sibling processes inside one pod. It
confirmed the arithmetic of the accounting. Measured over predicted memory was
0.999--1.000 at the median for every staging, and tmpfs was charged at exactly the input
size. A pod with one container has one memory cgroup, so this check could not show
where the kernel charges pages when a pod has several containers.

The second check (rt-k8spod-1) used real pods. We took 100 runs (48 from RightSizeMem, 44
from RightSizeCPU, and 8 init-tmpfs runs with inputs over 300\,MiB) and gave each of them
pods of its own, about 480 in all (\cref{tab:podprobe}). Every container sampled its own
cgroup's memory counters every 50 milliseconds.

\begin{table}[tbp]
\centering\small
\caption{The pods in the real-pod check (rt-k8spod-1).}
\label{tab:podprobe}
\begin{tabular}{@{}lp{0.62\linewidth}@{}}
\toprule
pod & what it measures \\
\midrule
Flyte task & the run as a real Flyte task: the rendered module in a fresh task pod, with its inputs downloaded by Flyte \\
Job, generous limit & the pod the manifest describes, with each container reading its own cgroup \\
Job at its label & does the run fit when the job container's limit is set to the label? \\
Job one cell lower & does the cell below the label fail? \\
init-tmpfs extras & the first release's flat 256\,Mi init container, and the job at its label minus the staged bytes \\
\bottomrule
\end{tabular}
\end{table}

\paragraph{The Flyte runtime.}
Before a task's own module is imported, the Flyte task container already holds 55.0\,MiB
of anonymous memory, with almost no variation across the 100 runs (p10--p90
54.9--55.0\,MiB). Its peak is a median 70\,MiB above the same run's peak as a plain Job
(p10 53, p90 242\,MiB), so running a workload inside the runtime's process costs slightly
more than the runtime does on its own. With the first release's 180\,MiB, the Flyte label matched the real container's cell
for only 23\% of runs and was one cell too high for 70\%. With 70\,MiB, it matches for
58\%, is one cell high for 33\% and one cell low for 9\%. The low ones are process pools,
where the container counts every worker process, and some DuckDB and polars runs. The
k8s-2 release uses 70\,MiB, which lowers 61\% of the Flyte labels on CPU nodes by one
cell.

\paragraph{Where tmpfs pages are charged.}
A tmpfs page is charged to the container that writes it. When the job downloads its own
inputs (boto3-tmpfs), the job container's shared memory grows by the staged size times
1.0001 at the median ($n=23$). When an init container stages them (init-tmpfs), the init
container's shared memory equals the staged size and the job container's stays at zero
($n=31$). The pages outlive the init container, however, and still count against the
pod. With its limit set to the label minus the staged bytes, the job was OOM-killed in 9
of 31 init-tmpfs pods, and with the full label 30 of 31 ran. Adding the staged bytes to the label is therefore correct; the pages are enforced
through the pod's limit, not the container's.

\paragraph{The first release's init containers.}
With the flat 256\,Mi init container of the first release, every staging up to 162\,MiB
ran and none of 254\,MiB or more did. Ten of those were confirmed OOM-killed while staging,
and six more failed. In that release, 982 of the 2,756 evaluated init-tmpfs Jobs staged
more than 200\,MiB, which means 7\% of all Job rows described pods that could never have
started. The k8s-2 release sizes the init container at 256\,MiB plus the staged bytes.

\paragraph{Labels as limits.}
When we set the job container's limit to its label, 98 of 100 runs succeeded
(\cref{tab:limits}). Both failures were process pools, where the label records the largest
single process but the cgroup counts all of them together. One cell below the label, 25
of the 97 runs that started were OOM-killed.

\begin{table}[tbp]
\centering\small
\caption{Job pods with the limit at the label cell and one cell below it (rt-k8spod-1).
``Other'' is a failure that was not an OOM kill.}
\label{tab:limits}
\begin{tabular}{@{}lrrrrrr@{}}
\toprule
& \multicolumn{3}{c}{at label} & \multicolumn{3}{c}{one cell below} \\
\cmidrule(lr){2-4}\cmidrule(l){5-7}
staging & ran & OOM & other & ran & OOM & other \\
\midrule
init-disk   & 22 & 1 & 0 & 17 & 5 & 1 \\
init-tmpfs  & 30 & 1 & 0 & 24 & 7 & 0 \\
boto3-disk  & 23 & 0 & 0 & 16 & 6 & 1 \\
boto3-tmpfs & 23 & 0 & 0 & 15 & 7 & 1 \\
\bottomrule
\end{tabular}
\end{table}

\paragraph{Labels are sometimes a cell too generous.}
We also compared each corpus label with the job container's measured peak of anonymous
plus shared memory, plus the same 5\% headroom (\cref{tab:cells}). Depending on the
staging, the label falls in the same cell for 52--68\% of runs and within one cell for
91--100\%. It is higher for 26--48\% of runs and lower for at most 6\%. Two things push
labels up. The first is a quirk of \code{ru\_maxrss}: on Linux, the value a parent reads
for its child through \code{wait4} includes the parent's own high-water mark at the moment
the child calls \code{exec}, because the kernel carries the old address space's peak into
the new one. The measuring parent in our corpora held about 158\,MiB, so even the smallest
workloads report at least that much. The second is memory-mapped input files, which RSS
counts but the kernel can reclaim; \code{numpy\_linalg} and \code{image\_batch} both map
their inputs. Neither effect makes a run fail at its label, and together they make about a third of
labels one cell larger than necessary. The same comparison explains why 72 of the 97 Job pods
that started one cell below their label still ran: for many of them the label is one cell high.
The simulator of \cref{sec:cost} replays the measured need, so the cheapest request it finds
inherits the same bias, and regret is measured against a slightly generous floor.

\begin{table}[tbp]
\centering\small
\caption{The corpus Job label compared with the job container's measured cgroup peak, in
grid cells.}
\label{tab:cells}
\begin{tabular}{@{}lrrrrr@{}}
\toprule
staging & $n$ & same cell & within one & higher & lower \\
\midrule
init-disk   & 23 & 65\% &  91\% & 30\% & 4\% \\
init-tmpfs  & 31 & 68\% &  97\% & 26\% & 6\% \\
boto3-disk  & 23 & 52\% &  96\% & 48\% & 0\% \\
boto3-tmpfs & 23 & 52\% & 100\% & 48\% & 0\% \\
\bottomrule
\end{tabular}
\end{table}
